\section{Introduction}

\subsection{The assignment}

The CS4110 final examination asks each group of three to build a mini-language analyzer
over real, manually transcribed speech collected in Yaound\'e, a city where everyday
conversation mixes French, English, Pidgin, Fulfulde and Ewondo with locally invented
slang. The graded work has five components: a corpus of 10--15 authentic statements; a
lexical specification with regular-expression token types and a custom lexical analyzer;
a context-free grammar reduced to LL(1) form with left-recursion removal, left factoring,
FIRST/FOLLOW sets and a parsing table; a working parser that accepts or rejects the
collected sentences; and this report together with a live demonstration.

\subsection{What we built}

Rather than a single script, we built \textbf{Francanglais Studio}: a complete compiler
front end -- scanner, lexer, grammar preparation, LL(1) table construction and a
table-driven parser -- wrapped in a web application so that every compiler artefact the
rubric asks to see is directly inspectable rather than buried in console output.

\begin{keypoint}[Design principle]
The analyzer recognises a \emph{documented subset} of collected speech. It never pretends
to understand every sentence, never silently drops an unknown word, and never silently
resolves a grammar ambiguity. Unknown vocabulary and parsing conflicts are both reported
explicitly, because those reports are the evidence this report is built on.
\end{keypoint}

\subsection{Scale of the artefact}

\begin{center}
\begin{tabular}{@{}lr@{}}
\toprule
\textbf{Artefact} & \textbf{Size} \\
\midrule
Lexicon entries & \LexiconTotal \\
\quad corpus-attested & \LexiconCorpus \\
\quad dictionary-derived & \LexiconDictionary \\
\quad multiword entries & \LexiconMultiword \\
Grammar terminals & \TerminalCount \\
Grammar nonterminals (as authored) & \NonterminalCount \\
Productions (authored / executable) & \ProductionCount\ / \ExecProductionCount \\
Transformation steps recorded & \LedgerSteps \\
Donor languages modelled & \OriginCount \\
Topic rules & \TopicCount \\
\bottomrule
\end{tabular}
\end{center}

Every number above, and every table in this report, is exported directly from the running
analyzer by \texttt{tools/export\_tables.py}. Nothing is transcribed by hand, so the report
cannot drift from the code it documents.

\subsection{How to read this report}

Sections~\ref{sec:data}--\ref{sec:results} follow the marking scheme in order: data
collection, lexical analysis, syntactic analysis, implementation, and the accept/reject
results. Section~\ref{sec:complexity} answers the required discussion question on why
Yaound\'e communication is linguistically complex. Section~\ref{sec:architecture} documents
the surrounding software architecture. Appendices hold the full grammar listings, the
corpus-attested lexicon and screenshots of the working analyzer.
